\section{How many zeros distinguish two couplings?}\label{sec:couplinggap}

Fix a common ordered list $\cX$ of $n$ distinct characters. Set
$p_A(y)=\det(A+\diag y)$ and $p_B(y)=\det(B+\diag y)$ for $y\in\CC^n$.
The coefficient of $\prod_{i\notin I}y_i$ is the principal minor on $I$.
Write $M_{A\to B}(I)$ for the one-sided discrepancy of the two coupled
divisors on a height interval $I$.

\begin{theorem}[Quantitative coupling separation]\label{thm:couplinggap}
Assume Dirichlet GRH. If $p_A\ne p_B$, there exist $j\in\{1,\ldots,n+1\}$,
$c>0$, and $T_0$ such that for every $T\ge T_0$,
\[
 \min\{M_{A\to B}([T/2,3T]),\,M_{B\to A}([T/2,3T])\}
       \ \ge\ c\,T/(\log T)^j .
\]
If $r$ is the smallest size of an unequal principal minor, one can take
$j\le\min\{n+1,2r\}$. For Laplacians of graphs with nonnegative
weights this gives $j\le\min\{n+1,4\}$. The constants may depend on the two matrices and the
characters. Consequently $\liminf_{T\to\infty}
M_{A\to B}((0,T])(\log T)^{n+1}/T=0$ forces equality of all principal
minors. For these graph Laplacians the fourth power suffices in every
dimension and the aligned matrices are equal.
\end{theorem}

\begin{proof}
For a formal variable $u$ and independent $y=(y_1,\ldots,y_n)$ write
\begin{equation}\label{eq:coupling-logexp}
 \log\frac{p_A(u\mathbf1+y)}{p_B(u\mathbf1+y)}
 =\sum_{k\ge1}h_k(y)u^{-k},\qquad
 h_k(y)=\frac{(-1)^{k+1}}k
 \tr\big[(A+\diag y)^k-(B+\diag y)^k\big].
\end{equation}
This is the power series with zero constant term; it converges for all
sufficiently large $|u|$, locally uniformly in $y$, with termwise
derivatives. Let $k_0$ be its first nonzero index. The degree in $u$ of
$p_A(u\mathbf1+y)-p_B(u\mathbf1+y)$ is some $\ell\in[0,n-1]$ over
$\CC[y]$, since setting $u=0$ gives the nonzero polynomial $p_A-p_B$.
The denominator is monic of degree $n$, so $k_0=n-\ell\le n$. If
$h_{k_0}$ is a nonzero constant $d$, replacing $y$ by $y+v\mathbf1$
is the same as replacing $u$ by $u+v$. Coefficients at the next power give
$h_{k_0+1}(y+v\mathbf1)=h_{k_0+1}(y)-k_0dv$.
The first nonconstant coefficient $h_j$ therefore exists and has
$j\le n+1$.

To see the finer bound, set $v_i=(u+y_i)^{-1}$ and expand
\[
 \frac{p_A(u\mathbf1+y)}{\prod_i(u+y_i)}
   =1+\sum_{I\ne\varnothing}\det A[I]\prod_{i\in I}v_i.
\]
In the logarithm of the ratio for $A,B$, all monomials in $v$ of total
degree less than $r$ vanish, and its homogeneous degree-$r$ part is
$\sum_{|I|=r}(\det A[I]-\det B[I])\prod_{i\in I}v_i$. Indeed,
dividing the difference by the second normalized polynomial and then
taking $\log(1+x)$ leave the first nonzero homogeneous part unchanged.
At order $u^{-2r}$, a monomial of $v$-degree $d$ yields a polynomial
homogeneous in $y$ of degree $2r-d$ (or contributes nothing). For
a fixed $r$-set $I$, the coefficient of $\prod_{i\in I} y_i$ in
$h_{2r}$ is therefore
$(-1)^r(\det A[I]-\det B[I])$, since each of the corresponding $r$
factors must supply its linear $y_i$ term. Some such coefficient
is nonzero. For Laplacians the diagonal entries give the order-one
minors, and $\det A[\{i,k\}]=A_{ii}A_{kk}-|A_{ik}|^2$ together
with $A_{ik}\le0$ determines every edge weight. Distinct
nonnegative weighted graph Laplacians therefore have $r\le2$.

Put $D_i(s)=iL'(s,\chi_i)/L(s,\chi_i)$ and
$c_i=(i/2)\log(q_i/\pi)$. The squarefree-prime argument in
Lemma~\ref{indep} gives algebraic independence of the $D_i$ over $\CC$.
It follows that $h_j(c+D(s))$ is a nonconstant Dirichlet series.

To detect this coefficient in the zeros, write
$s=2+it$, $U=(i/2)\log(s/2)$, and
$g_i(-i(s-1/2))=U+y_i(s)$.
Stirling's formula for the digamma function and its derivative gives
$y_i=c_i+D_i+O(t^{-1})$ and
$y_i'=D_i'+O(t^{-1})$ as $t\to+\infty$, where primes here mean
$d/ds$. Both $y$ and $y'$ range over bounded sets.
The product of the input functions cancels in $P_A(s)/P_B(s)$,
where $P_A(s)=\PP_{\cX,A}(-i(s-1/2))$. Differentiate
\eqref{eq:coupling-logexp} at $(U,y(s))$. Lower coefficients are
constant, $U'=O(t^{-1})$, and the differentiated tails converge
uniformly because the bounded $y$ lie in a compact set. Thus
\[
 \frac{P_A'}{P_A}-\frac{P_B'}{P_B}
  = U^{-j}\frac{d}{ds}h_j(c+D(s))
     +O((\log t)^{-j-1})+O(t^{-1}).
\]
The right side is a logarithmic derivative with respect to $s$, as in
\eqref{poiss}. In particular $U=(i/2)\log t+O(1)$ gives
\begin{align}
 H_A(t)-H_B(t)&=(\log t)^{-j} f(t)+O((\log t)^{-j-1}),\label{eq:gapsig}\\
 f(t)&=\Re\left[ (2/i)^j\frac{d}{ds}h_j(c+D(s))\bigg|_{s=2+it}\right].\nonumber
\end{align}
Here $H_A=\Re(P_A'/P_A)$, and similarly for $B$.

The complex series inside the real part has only frequencies $-\log m$,
$m\ge2$, with at least one nonzero coefficient. Its differentiated
Dirichlet series is absolutely convergent at $2$ by absolute convergence
of the undifferentiated series on every half-plane $\Re s>1$.
Taking the real part cannot cancel a nonzero coefficient against its
conjugate, which occurs at the opposite frequency. Hence $f$ has mean
zero and positive mean square. More explicitly, uniform approximation
by finite exponential sums and integration of their frequencies give
$T^{-1}\int_T^{2T} f\to0$ and
$T^{-1}\int_T^{2T} f^2\to m_2>0$; the latter limit is one half the sum
of the squared moduli of the complex coefficients. Since
$|f|\ge f^2/\|f\|_\infty$ and
$2f_\pm=|f|\pm f$, each of $\int_T^{2T} f_\pm$ is $\gg T$.
Taking positive and negative parts of \eqref{eq:gapsig} therefore gives
\[
 \int_T^{2T}(H_A-H_B)_\pm\,dt \gg T/(\log T)^j .
\]

Finally let $\nu=Z_{\cX,A}-Z_{\cX,B}$ as a signed counting measure,
and $K_b(v)=b/(b^2+v^2)$, $b=3/2$.
The positivity in \eqref{poiss} implies
$(H_A-H_B)_\pm\le K_b*\nu_\pm$.
The local zero estimate following \eqref{cline} bounds each of the
two measures on $[v,v+1]$ by $O(\log(2+|v|))$.
Every atom inside $[T/2,3T]$ contributes at most $\pi$ to the
integral over $[T,2T]$. The contribution of all outside atoms is
$O(\log T)$: those with $|v|\le4T$ contribute $O(1/T)$ each by
distance at least $T/2$, and those farther away contribute $O(T/v^2)$
each. Combining the estimates proves the window bound. Replacing $T$ by
$T/3$ gives the consequence on $(0,T]$.
\end{proof}
